\documentclass[11pt,a4paper]{article}

% ── Paquetes ──────────────────────────────────────────────────────────────
\usepackage[utf8]{inputenc}
\usepackage[T1]{fontenc}
\usepackage[spanish,es-tabla]{babel}
\usepackage{amsmath,amssymb,amsthm}
\usepackage{mathtools}
\usepackage{geometry}
\usepackage{enumitem}
\usepackage{booktabs}
\usepackage{array}
\usepackage{tikz}
\usetikzlibrary{automata,positioning,arrows.meta,babel}
\usepackage{hyperref}
\usepackage{xcolor}
\usepackage{framed}
\usepackage{tcolorbox}
\tcbuselibrary{theorems,breakable}

\geometry{margin=2.5cm}

% ── Entornos ──────────────────────────────────────────────────────────────
\theoremstyle{definition}
\newtheorem{definicion}{Definición}[section]
\newtheorem{teorema}{Teorema}[section]
\newtheorem{proposicion}{Proposición}[section]
\newtheorem{corolario}{Corolario}[section]
\newtheorem{lema}{Lema}[section]

\theoremstyle{remark}
\newtheorem{observacion}{Observación}[section]

\newtcolorbox{ejemplo}[1][]{%
  colback=blue!5!white,
  colframe=blue!60!black,
  fonttitle=\bfseries,
  title=Ejemplo: #1,
  breakable
}

\newtcolorbox{ejercicio}[1][]{%
  colback=green!5!white,
  colframe=green!50!black,
  fonttitle=\bfseries,
  title=Ejercicio propuesto: #1,
  breakable
}

\newtcolorbox{cajanota}[1][]{%
  colback=yellow!8!white,
  colframe=orange!70!black,
  fonttitle=\bfseries,
  title=Nota importante: #1,
  breakable
}

% ── Comandos útiles ───────────────────────────────────────────────────────
\newcommand{\R}{\mathbb{R}}
\newcommand{\N}{\mathbb{N}}
\newcommand{\E}{\mathbb{E}}
\newcommand{\Var}{\mathrm{Var}}
\newcommand{\Cov}{\mathrm{Cov}}
\newcommand{\Prob}{\mathbb{P}}
\newcommand{\e}{\mathrm{e}}
\DeclareMathOperator{\diag}{diag}

% ══════════════════════════════════════════════════════════════════════════
\title{%
  \textbf{Análisis de Estado Estable de Cadenas de Markov\\
  de Tiempo Continuo}\\[6pt]
  \large Lectura para curso de pregrado
}
\author{Notas de clase}
\date{2026}

\begin{document}
\maketitle

\tableofcontents
\newpage

% ======================================================================
%  SECCIÓN 1 — MOTIVACIÓN
% ======================================================================
\section{Motivación: ¿qué describe el estado estable?}

En la lectura anterior sobre análisis transitorio vimos cómo calcular la
distribución $\boldsymbol{p}(t) = \boldsymbol{\alpha}\,\e^{Qt}$ para un
instante particular $t$.  Observamos que, en todos los ejemplos, cuando
$t$ crecía, las probabilidades convergían a valores fijos
independientes de la distribución inicial $\boldsymbol{\alpha}$.  Esa
distribución límite es la \textbf{distribución de estado estable}
(o \emph{estacionaria}).

¿Por qué nos importa?

\begin{itemize}[nosep]
  \item Un servidor web lleva meses operando.  Las condiciones iniciales
        ya no son relevantes; queremos saber la fracción del tiempo que
        está saturado.
  \item Una clínica veterinaria quiere dimensionar su sala de espera.
        Le interesa el número promedio de mascotas en cola \emph{a largo
        plazo}, no en los primeros minutos del día.
  \item Un ingeniero de redes necesita calcular la tasa promedio de
        paquetes perdidos cuando un enlace ha estado operando por tiempo
        indefinido.
  \item Un gerente de banco quiere saber qué porcentaje del tiempo sus
        cajeros están ociosos, para decidir si contrata uno más.
\end{itemize}

En todos estos casos, el sistema ha ``olvidado'' sus condiciones
iniciales y opera en un régimen permanente.  El análisis de estado
estable nos da las herramientas para responder estas preguntas sin
necesidad de calcular la exponencial de matriz.

% ======================================================================
%  SECCIÓN 2 — EXISTENCIA Y UNICIDAD
% ======================================================================
\section{Existencia y unicidad de la distribución estacionaria}

\subsection{De lo transitorio a lo estable: el paso al límite}

Recordemos que la distribución transitoria satisface
\[
  \boldsymbol{p}(t) = \boldsymbol{\alpha}\,\e^{Qt}.
\]
Si definimos $\boldsymbol{\pi} = \lim_{t\to\infty}\boldsymbol{p}(t)$
(asumiendo que el límite existe), entonces al derivar
$\boldsymbol{p}'(t) = \boldsymbol{p}(t)\,Q$ y tomar $t\to\infty$:
\[
  \lim_{t\to\infty}\boldsymbol{p}'(t) = \boldsymbol{0}
  \quad\Longrightarrow\quad
  \boldsymbol{\pi}\,Q = \boldsymbol{0}.
\]

Esta es la intuición: en estado estable la distribución ya no cambia,
así que su derivada temporal es cero.

\subsection{Condiciones formales}

\begin{teorema}[Existencia y unicidad]
\label{teo:existencia}
Sea $\{X(t)\}$ una CTMC con espacio de estados finito
$\mathcal{S} = \{0,1,\dots,n-1\}$ y generadora $Q$.  Si la cadena es
\textbf{irreducible} (todo estado se puede alcanzar desde cualquier
otro), entonces:
\begin{enumerate}[nosep]
  \item Existe un único vector $\boldsymbol{\pi} = (\pi_0,\pi_1,\dots,\pi_{n-1})$ tal que
    \[
      \boldsymbol{\pi}\,Q = \boldsymbol{0}, \qquad
      \sum_{j=0}^{n-1}\pi_j = 1, \qquad
      \pi_j > 0 \;\;\forall j.
    \]
  \item Para cualquier distribución inicial $\boldsymbol{\alpha}$,
    \[
      \lim_{t\to\infty}\boldsymbol{\alpha}\,\e^{Qt} = \boldsymbol{\pi}.
    \]
  \item $\pi_j$ tiene dos interpretaciones equivalentes:
    \begin{itemize}[nosep]
      \item Probabilidad a largo plazo de estar en el estado $j$.
      \item Fracción del tiempo que el sistema pasa en el estado $j$.
    \end{itemize}
\end{enumerate}
\end{teorema}

\begin{cajanota}[¿Por qué se necesita irreducibilidad?]
Si la cadena no es irreducible, puede haber múltiples clases
comunicantes, cada una con su propia distribución estacionaria.  En ese
caso, la distribución límite \emph{sí} depende de $\boldsymbol{\alpha}$:
el sistema queda ``atrapado'' en la clase a la que llegue.  Para esta
lectura asumiremos siempre irreducibilidad.
\end{cajanota}

\subsection{Deducción algebraica del sistema de ecuaciones}

La ecuación $\boldsymbol{\pi}Q = \boldsymbol{0}$ se escribe componente
a componente como
\[
  \sum_{i=0}^{n-1}\pi_i\,q_{ij} = 0, \qquad j = 0,1,\dots,n-1.
\]
Separando el término diagonal:
\[
  \pi_j\,q_{jj} + \sum_{i\ne j}\pi_i\,q_{ij} = 0
  \quad\Longrightarrow\quad
  \pi_j\,|q_{jj}| = \sum_{i\ne j}\pi_i\,q_{ij}.
\]
Dado que $|q_{jj}| = \sum_{k\ne j}q_{jk}$, esta ecuación dice:

\begin{center}
\fbox{\parbox{0.8\textwidth}{
\centering
\textbf{Tasa total de salida del estado $j$}
= \textbf{Tasa total de entrada al estado $j$}.
}}
\end{center}

Esto se conoce como las \textbf{ecuaciones de balance global}.

% ======================================================================
%  SECCIÓN 3 — ECUACIONES DE BALANCE
% ======================================================================
\section{Ecuaciones de balance}

\subsection{Balance global}

Para cada estado $j$:
\[
  \underbrace{\pi_j \sum_{k\ne j} q_{jk}}_{\text{flujo que sale de }j}
  \;=\;
  \underbrace{\sum_{i\ne j}\pi_i\,q_{ij}}_{\text{flujo que entra a }j}.
\]

Son $n$ ecuaciones, pero como las filas de $Q$ suman cero, una ecuación
es combinación lineal de las demás.  Se descarta una y se reemplaza por
la condición de normalización $\sum_j \pi_j = 1$.

\subsection{Balance detallado (\textit{detailed balance})}

En algunos sistemas (particularmente los reversibles), se cumple una
condición más fuerte:
\[
  \pi_i\,q_{ij} = \pi_j\,q_{ji}, \qquad \forall\; i \ne j.
\]
Esto dice que el flujo de $i$ a $j$ es igual al flujo de $j$ a $i$
para \emph{cada par} de estados.  Si se cumple balance detallado,
entonces automáticamente se cumple balance global (sumar sobre $i$).

\begin{cajanota}[Utilidad del balance detallado]
El balance detallado reduce enormemente la cantidad de trabajo.  En
lugar de resolver un sistema de ecuaciones, se obtienen relaciones
directas entre las $\pi_j$.  Esto es particularmente útil en cadenas
de nacimiento y muerte, donde las únicas transiciones son entre estados
vecinos.
\end{cajanota}

\subsection{Cadenas de nacimiento y muerte}

Para una cadena con estados $\{0,1,\dots,K\}$ donde solo hay
transiciones entre vecinos (tasas $\lambda_j$ de $j$ a $j+1$ y $\mu_j$
de $j$ a $j-1$), el balance detallado da:
\[
  \pi_j\,\lambda_j = \pi_{j+1}\,\mu_{j+1}, \qquad j=0,1,\dots,K-1.
\]
Despejando recursivamente:
\begin{align*}
  \pi_1 &= \frac{\lambda_0}{\mu_1}\,\pi_0, \\
  \pi_2 &= \frac{\lambda_1}{\mu_2}\,\pi_1
          = \frac{\lambda_0\lambda_1}{\mu_1\mu_2}\,\pi_0, \\
  &\;\;\vdots \\
  \pi_j &= \frac{\lambda_0\lambda_1\cdots\lambda_{j-1}}
                 {\mu_1\mu_2\cdots\mu_j}\,\pi_0
         = \pi_0\prod_{k=0}^{j-1}\frac{\lambda_k}{\mu_{k+1}}.
\end{align*}
Finalmente, $\pi_0$ se obtiene de $\sum_{j=0}^K \pi_j = 1$.

% ======================================================================
%  SECCIÓN 4 — MÉTODOS DE SOLUCIÓN
% ======================================================================
\section{Métodos para resolver \texorpdfstring{$\boldsymbol{\pi}Q=\boldsymbol{0}$}{πQ=0}}

\subsection{Método 1: sustitución directa (sistemas pequeños)}

Para 2 o 3 estados, se escribe el sistema a mano, se descarta una
ecuación redundante, se sustituye la normalización, y se resuelve
por sustitución o eliminación.

\subsection{Método 2: balance detallado (si aplica)}

Si la cadena es de nacimiento y muerte (o más generalmente, reversible),
se usan las relaciones $\pi_i q_{ij} = \pi_j q_{ji}$ para expresar
todo en función de $\pi_0$, y luego se normaliza.

\subsection{Método 3: forma matricial}

Se reemplaza una columna de $Q^\top$ por un vector de unos (para
incorporar la normalización) y se resuelve el sistema lineal resultante.
Concretamente, sea $\tilde{Q}$ la matriz que se obtiene de $Q^\top$ al
reemplazar la última fila por $(1,1,\dots,1)$.  Entonces
\[
  \tilde{Q}\,\boldsymbol{\pi}^\top = \begin{pmatrix}0\\0\\\vdots\\0\\1\end{pmatrix}
  \quad\Longrightarrow\quad
  \boldsymbol{\pi}^\top = \tilde{Q}^{-1}\begin{pmatrix}0\\\vdots\\0\\1\end{pmatrix}.
\]

\subsection{Método 4: a partir de la cadena embebida}

La cadena embebida de tiempo discreto tiene matriz de transición
$\hat{P}$ con entradas $\hat{p}_{ij} = q_{ij}/|q_{ii}|$ para $i\ne j$.
Si $\hat{\boldsymbol{\pi}}$ es la distribución estacionaria de
$\hat{P}$, entonces la distribución estacionaria de la CTMC es
\[
  \pi_j = \frac{\hat{\pi}_j / |q_{jj}|}
               {\sum_k \hat{\pi}_k / |q_{kk}|}.
\]
La intuición: la cadena embebida indica con qué frecuencia se
\emph{visita} cada estado, pero el tiempo de permanencia en cada visita
es $1/|q_{jj}|$.  La distribución de tiempo continuo pondera por ambos
factores.

% ======================================================================
%  SECCIÓN 5 — MÉTRICAS DE DESEMPEÑO
% ======================================================================
\section{Métricas de desempeño en estado estable}

Una vez que tenemos $\boldsymbol{\pi}$, podemos calcular cualquier
cantidad que dependa del estado del sistema ``a largo plazo''.

\subsection{Valor esperado en estado estable}

Para una función $f:\mathcal{S}\to\R$:
\[
  \E_\pi[f(X)] = \sum_{j\in\mathcal{S}} f(j)\,\pi_j.
\]

\subsection{Varianza en estado estable}

\[
  \Var_\pi[f(X)] = \sum_{j\in\mathcal{S}} f(j)^2\,\pi_j
  - \Bigl(\sum_{j\in\mathcal{S}} f(j)\,\pi_j\Bigr)^{\!2}
  = \E_\pi[f(X)^2] - \bigl(\E_\pi[f(X)]\bigr)^2.
\]

\begin{cajanota}[Deducción detallada de la varianza]
La varianza se deduce igual que en el caso transitorio.  Sea
$Y = f(X)$ donde $X$ tiene distribución $\boldsymbol{\pi}$.
\begin{align*}
  \Var[Y] &= \E[(Y-\E[Y])^2] = \E[Y^2 - 2Y\E[Y] + (\E[Y])^2] \\
  &= \E[Y^2] - 2\E[Y]\E[Y] + (\E[Y])^2 = \E[Y^2] - (\E[Y])^2.
\end{align*}
Expandiendo:
\[
  \E[Y^2] = \sum_j f(j)^2\,\pi_j, \qquad
  (\E[Y])^2 = \Bigl(\sum_j f(j)\,\pi_j\Bigr)^{\!2}.
\]
No hay ningún paso especial de cadenas de Markov; es álgebra de
variables aleatorias discretas aplicada a la distribución $\pi$.
\end{cajanota}

\subsection{Tasas de flujo en estado estable}

La \textbf{tasa de flujo de $i$ a $j$} en estado estable es
$\pi_i\,q_{ij}$.  Esto permite calcular, por ejemplo:

\begin{itemize}[nosep]
  \item \textbf{Tasa efectiva de llegada a un sistema con pérdidas:}
    la tasa de entrada real (no la ofrecida) es
    $\lambda_{\text{ef}} = \lambda\sum_{j:\text{hay espacio}}\pi_j$.
  \item \textbf{Tasa de pérdida:} $\lambda_{\text{perd}} = \lambda\,\pi_K$
    si el estado $K$ es ``lleno''.
  \item \textbf{Throughput:} tasa de servicio completado,
    $\mu\sum_{j:\text{servidor ocupado}}\pi_j$.
\end{itemize}

\subsection{Métricas integrales: conexión con el tiempo}

Puesto que $\pi_j$ es la fracción del tiempo en estado $j$:
\[
  \E_\pi[f(X)] = \lim_{T\to\infty}\frac{1}{T}\int_0^T f(X(s))\,ds
  \quad\text{(con probabilidad 1).}
\]
Esto es el \textbf{teorema ergódico}.  La media temporal iguala la media
espacial.  Formalmente, la integral del lado derecho es una variable
aleatoria que converge casi seguramente al valor determinista del lado
izquierdo.

% ======================================================================
%  SECCIÓN 6 — EJEMPLO 1: DOS ESTADOS
% ======================================================================
\section{Ejemplo 1 --- Máquina con fallas (dos estados)}

\begin{ejemplo}[Máquina industrial que falla y se repara]
Una máquina en una planta de producción puede estar operativa (estado~0)
o descompuesta (estado~1).  Las fallas ocurren con tasa
$\lambda = 0.5$ fallas/hora y las reparaciones con tasa
$\mu = 1.5$ reparaciones/hora.  La máquina lleva operando meses y
queremos responder:

\begin{enumerate}[label=(\alph*)]
  \item ¿Qué fracción del tiempo está descompuesta?
  \item Si cada hora de inactividad cuesta \$200, ¿cuál es el costo
        promedio por hora?
  \item ¿Cuál es la varianza del costo por hora?
\end{enumerate}

\tcblower

\textbf{Paso 1: la generadora.}

\begin{center}
\begin{tikzpicture}[->,>=Stealth,node distance=3cm,
    every state/.style={thick,fill=blue!10}]
  \node[state] (0) {0};
  \node[state] (1) [right=of 0] {1};
  \path (0) edge[bend left] node[above] {$\lambda=0.5$} (1)
        (1) edge[bend left] node[below] {$\mu=1.5$} (0);
\end{tikzpicture}
\end{center}

\[
  Q = \begin{pmatrix} -0.5 & 0.5 \\ 1.5 & -1.5 \end{pmatrix}.
\]

\textbf{Paso 2: resolver $\boldsymbol{\pi}Q = \boldsymbol{0}$.}

El sistema es:
\[
  \begin{cases}
    -0.5\,\pi_0 + 1.5\,\pi_1 = 0, \\
    0.5\,\pi_0 - 1.5\,\pi_1 = 0.
  \end{cases}
\]
Ambas ecuaciones dicen lo mismo: $\pi_0 = 3\pi_1$.  Con la normalización
$\pi_0 + \pi_1 = 1$:
\[
  3\pi_1 + \pi_1 = 1 \quad\Longrightarrow\quad \pi_1 = \frac{1}{4},
  \qquad \pi_0 = \frac{3}{4}.
\]

\textbf{Verificación por balance detallado:}
$\pi_0\,q_{01} = \frac{3}{4}\times 0.5 = 0.375$ y
$\pi_1\,q_{10} = \frac{1}{4}\times 1.5 = 0.375$. \checkmark

\textbf{Verificación alternativa:} para un sistema de dos estados con
tasas $\lambda$ y $\mu$, la fórmula general es
\[
  \pi_0 = \frac{\mu}{\lambda+\mu}, \qquad
  \pi_1 = \frac{\lambda}{\lambda+\mu}.
\]
Con $\lambda=0.5$, $\mu=1.5$: $\pi_0 = 1.5/2 = 3/4$,
$\pi_1 = 0.5/2 = 1/4$. \checkmark

\textbf{Paso 3: respuestas.}

(a) La máquina está descompuesta el $25\%$ del tiempo.

(b) Definimos $f(0) = 0$ (operativa, sin costo) y $f(1) = 200$
(descompuesta, \$200/hora):
\[
  \E_\pi[f(X)] = 0 \times \frac{3}{4} + 200 \times \frac{1}{4}
  = \$50 \text{ por hora}.
\]

(c) Varianza del costo:
\[
  \E_\pi[f(X)^2] = 0^2 \times \frac{3}{4} + 200^2 \times \frac{1}{4}
  = 10{,}000.
\]
\[
  \Var_\pi[f(X)] = 10{,}000 - 50^2 = 10{,}000 - 2{,}500 = 7{,}500.
\]
\[
  \sigma = \sqrt{7{,}500} \approx \$86.60 \text{ por hora}.
\]

La desviación estándar es alta relativa a la media, lo cual refleja la
naturaleza ``todo o nada'' del costo: o se pagan \$0 o \$200, con nada
intermedio.
\end{ejemplo}

% ======================================================================
%  SECCIÓN 7 — EJEMPLO 2: COLA M/M/1/2
% ======================================================================
\section{Ejemplo 2 --- Cola con capacidad limitada ($M/M/1/2$)}

\begin{ejemplo}[Fotocopiadora con espera limitada]
Una oficina tiene una fotocopiadora.  Los empleados llegan a usarla
con tasa $\lambda = 6$ personas/hora.  Cada trabajo tarda en promedio
8 minutos, así que $\mu = 60/8 = 7.5$ trabajos/hora.  Solo cabe una
persona esperando (además de la que está usando la máquina), así que
la capacidad es $K=2$.  Si alguien llega y hay 2 personas, se va a
otra planta.

La fotocopiadora lleva meses en servicio.  Se pide:

\begin{enumerate}[label=(\alph*)]
  \item Distribución estacionaria.
  \item Número promedio de personas en el sistema y su varianza.
  \item Fracción de empleados que se van sin fotocopiar (probabilidad de
        bloqueo).
  \item Tasa efectiva de uso de la fotocopiadora (throughput).
\end{enumerate}

\tcblower

\textbf{Paso 1: generadora.}

\begin{center}
\begin{tikzpicture}[->,>=Stealth,node distance=2.5cm,
    every state/.style={thick,fill=blue!10}]
  \node[state] (0) {0};
  \node[state] (1) [right=of 0] {1};
  \node[state] (2) [right=of 1] {2};
  \path (0) edge[bend left] node[above] {$6$} (1)
        (1) edge[bend left] node[below] {$7.5$} (0)
        (1) edge[bend left] node[above] {$6$} (2)
        (2) edge[bend left] node[below] {$7.5$} (1);
\end{tikzpicture}
\end{center}

\[
  Q = \begin{pmatrix}
    -6 & 6 & 0 \\
    7.5 & -13.5 & 6 \\
    0 & 7.5 & -7.5
  \end{pmatrix}.
\]

\textbf{Paso 2: resolver por balance detallado.}

Es una cadena de nacimiento y muerte con $\lambda_j = 6$ y $\mu_j = 7.5$
para todo $j$.  Definimos $\rho = \lambda/\mu = 6/7.5 = 0.8$.

Balance detallado entre estados vecinos:
\begin{align*}
  \pi_0 \cdot 6 &= \pi_1 \cdot 7.5
  \quad\Longrightarrow\quad \pi_1 = \rho\,\pi_0 = 0.8\,\pi_0, \\
  \pi_1 \cdot 6 &= \pi_2 \cdot 7.5
  \quad\Longrightarrow\quad \pi_2 = \rho\,\pi_1 = \rho^2\,\pi_0 = 0.64\,\pi_0.
\end{align*}
En general $\pi_j = \rho^j\,\pi_0$.  Normalización:
\[
  \pi_0(1 + \rho + \rho^2) = 1
  \quad\Longrightarrow\quad
  \pi_0 = \frac{1}{1+0.8+0.64} = \frac{1}{2.44} \approx 0.4098.
\]

\textbf{Fórmula general para $M/M/1/K$ con $\rho\ne 1$:}
\[
  \pi_0 = \frac{1-\rho}{1-\rho^{K+1}}, \qquad
  \pi_j = \rho^j\,\pi_0, \quad j=0,1,\dots,K.
\]
Verifiquemos: $\pi_0 = \frac{1-0.8}{1-0.8^3} = \frac{0.2}{1-0.512}
= \frac{0.2}{0.488} \approx 0.4098$. \checkmark

Distribución completa:
\[
  \boldsymbol{\pi} \approx (0.4098,\; 0.3279,\; 0.2623).
\]

\textbf{Paso 3: valor esperado.}

Sea $N$ el número de personas en el sistema, $f(j)=j$:
\[
  \E_\pi[N] = 0\times 0.4098 + 1\times 0.3279 + 2\times 0.2623
  = 0.3279 + 0.5246 = 0.8525.
\]

Para la fórmula cerrada, derivemos la expresión general.  Con
$\pi_j = \rho^j\pi_0$:
\begin{align*}
  \E[N] &= \sum_{j=0}^{K} j\,\rho^j\,\pi_0
  = \pi_0\sum_{j=1}^{K}j\,\rho^j
  = \pi_0\,\rho\frac{d}{d\rho}\sum_{j=0}^{K}\rho^j
  = \pi_0\,\rho\frac{d}{d\rho}\Bigl(\frac{1-\rho^{K+1}}{1-\rho}\Bigr).
\end{align*}
La derivada es:
\[
  \frac{d}{d\rho}\frac{1-\rho^{K+1}}{1-\rho}
  = \frac{-(K+1)\rho^K(1-\rho)+(1-\rho^{K+1})}{(1-\rho)^2}.
\]
Para $K=2$, $\rho=0.8$:
\[
  \frac{d}{d\rho}\frac{1-\rho^3}{1-\rho}\Bigg|_{\rho=0.8}
  = \frac{-3(0.64)(0.2)+0.488}{0.04}
  = \frac{-0.384+0.488}{0.04}
  = \frac{0.104}{0.04} = 2.60.
\]
\[
  \E[N] = 0.4098 \times 0.8 \times 2.60 = 0.8524 \approx 0.8525. \;\checkmark
\]

\textbf{Paso 4: varianza.}

Segundo momento:
\[
  \E[N^2] = 0^2\times 0.4098 + 1^2\times 0.3279 + 2^2\times 0.2623
  = 0.3279 + 1.0492 = 1.3771.
\]
\[
  \Var[N] = 1.3771 - (0.8525)^2 = 1.3771 - 0.7268 = 0.6503.
\]
\[
  \sigma_N = \sqrt{0.6503} \approx 0.8064.
\]

\textbf{Paso 5: probabilidad de bloqueo y throughput.}

Probabilidad de bloqueo (un cliente llega y el sistema está lleno):
\[
  P_{\text{bloqueo}} = \pi_2 \approx 0.2623.
\]
Es decir, aproximadamente el $26.23\%$ de los empleados se van sin
fotocopiar.

Tasa efectiva de llegada:
\[
  \lambda_{\text{ef}} = \lambda(1-\pi_2) = 6(1-0.2623) = 6\times 0.7377
  \approx 4.426 \text{ personas/hora}.
\]

Throughput (tasa de trabajos completados, que en estado estable iguala
a la tasa efectiva de llegada):
\[
  \text{Throughput} = \mu\,(1-\pi_0) = 7.5\times 0.5902 \approx 4.426. \;\checkmark
\]
Ambas formas dan el mismo resultado, como debe ser en estado estable.
\end{ejemplo}

% ======================================================================
%  SECCIÓN 8 — EJEMPLO 3: TRES ESTADOS, NO NACIMIENTO-MUERTE
% ======================================================================
\section{Ejemplo 3 --- Sistema que no es de nacimiento y muerte}

\begin{ejemplo}[Servidor web con modo degradado]
Un servidor web puede estar en tres estados: funcionando a plena
capacidad (estado~0), funcionando en modo degradado (estado~1), o
caído (estado~2).  Las transiciones son:

\begin{center}
\begin{tikzpicture}[->,>=Stealth,node distance=3cm,
    every state/.style={thick,fill=blue!10}]
  \node[state] (0) {0};
  \node[state] (1) [right=of 0] {1};
  \node[state] (2) [below right=1.5cm and 1.5cm of 0] {2};
  \path (0) edge[bend left=15] node[above] {$2$} (1)
        (0) edge node[left] {$0.5$} (2)
        (1) edge[bend left=15] node[below] {$3$} (0)
        (1) edge node[right] {$1$} (2)
        (2) edge node[below left] {$4$} (0);
\end{tikzpicture}
\end{center}

Note que desde el estado caído (2) solo se puede ir a plena capacidad
(0), no a degradado.  Esto \emph{no} es una cadena de nacimiento y
muerte (hay transición $0\to 2$ saltándose el estado~1).

El servidor lleva operando tiempo indefinido.  Si en estado~0 procesa
100 solicitudes/seg, en estado~1 procesa 40 solicitudes/seg, y en
estado~2 procesa 0, calcular el throughput promedio y su varianza.

\tcblower

\textbf{Paso 1: generadora.}

\[
  Q = \begin{pmatrix}
    -2.5 & 2 & 0.5 \\
    3 & -4 & 1 \\
    4 & 0 & -4
  \end{pmatrix}.
\]

\textbf{Paso 2: resolver $\boldsymbol{\pi}Q=\boldsymbol{0}$ por
sustitución.}

Las ecuaciones de balance global para cada estado son:

\emph{Estado 0 (flujo de entrada = flujo de salida):}
\[
  3\pi_1 + 4\pi_2 = 2.5\,\pi_0. \tag{E0}
\]

\emph{Estado 1:}
\[
  2\pi_0 = 4\,\pi_1. \tag{E1}
\]

\emph{Estado 2:}
\[
  0.5\,\pi_0 + 1\,\pi_1 = 4\,\pi_2. \tag{E2}
\]

De (E1): $\pi_1 = \pi_0/2$.

Sustituyendo en (E2): $0.5\pi_0 + 0.5\pi_0 = 4\pi_2$, es decir
$\pi_0 = 4\pi_2$, así que $\pi_2 = \pi_0/4$.

Verificamos en (E0): $3(\pi_0/2)+4(\pi_0/4) = 1.5\pi_0+\pi_0 = 2.5\pi_0$. \checkmark

Normalización:
\[
  \pi_0 + \frac{\pi_0}{2} + \frac{\pi_0}{4} = 1
  \quad\Longrightarrow\quad
  \frac{7}{4}\pi_0 = 1
  \quad\Longrightarrow\quad
  \pi_0 = \frac{4}{7}.
\]
\[
  \boldsymbol{\pi} = \Bigl(\frac{4}{7},\;\frac{2}{7},\;\frac{1}{7}\Bigr)
  \approx (0.5714,\;0.2857,\;0.1429).
\]

\textbf{Paso 3: throughput promedio.}

Sea $f(0)=100$, $f(1)=40$, $f(2)=0$ solicitudes/seg:
\[
  \E_\pi[f(X)] = 100\times\frac{4}{7} + 40\times\frac{2}{7}
  + 0\times\frac{1}{7}
  = \frac{400+80}{7} = \frac{480}{7} \approx 68.57 \text{ sol/seg}.
\]

\textbf{Paso 4: varianza del throughput.}
\[
  \E_\pi[f(X)^2] = 100^2\times\frac{4}{7} + 40^2\times\frac{2}{7}
  + 0^2\times\frac{1}{7}
  = \frac{40{,}000+3{,}200}{7} = \frac{43{,}200}{7} \approx 6{,}171.43.
\]
\[
  \Var_\pi[f(X)] = 6{,}171.43 - (68.57)^2 = 6{,}171.43 - 4{,}701.84
  = 1{,}469.59.
\]
\[
  \sigma \approx 38.34 \text{ solicitudes/seg}.
\]

\textbf{Interpretación:} el servidor procesa en promedio 68.6 sol/seg
a largo plazo, pero con una desviación de 38.3 sol/seg, reflejando que
oscila entre 100 y 0 frecuentemente.  El coeficiente de variación es
$38.3/68.6 \approx 56\%$, bastante alto.
\end{ejemplo}

% ======================================================================
%  SECCIÓN 9 — EJEMPLO CON DOS VARIABLES
% ======================================================================
\section{Ejemplo 4 --- Dos componentes con dependencia}

\begin{ejemplo}[Sistema de potencia con generador y respaldo]
Una estación tiene un generador principal (G) y un generador de respaldo
(R).  El generador G falla con tasa $\alpha=1$/día y se repara con
tasa $\beta=4$/día.  El respaldo R falla con tasa $\gamma=0.5$/día y
se repara con tasa $\delta=2$/día.

Aquí viene la dependencia: \textbf{solo hay un técnico de reparación}.
Si ambos generadores fallan simultáneamente, el técnico prioriza a G
(repara G primero; R espera).

El estado es $(G,R)$ donde $G,R\in\{0,1\}$ ($0$=funciona, $1$=fallado).
Los estados son: $(0,0)$, $(1,0)$, $(0,1)$, $(1,1)$, que numeramos
$0,1,2,3$.

\tcblower

\textbf{Paso 1: generadora.}

Las transiciones:
\begin{itemize}[nosep]
  \item $(0,0)\to(1,0)$: G falla, tasa $\alpha=1$.
  \item $(0,0)\to(0,1)$: R falla, tasa $\gamma=0.5$.
  \item $(1,0)\to(0,0)$: G se repara (técnico libre), tasa $\beta=4$.
  \item $(1,0)\to(1,1)$: R también falla, tasa $\gamma=0.5$.
  \item $(0,1)\to(0,0)$: R se repara (técnico libre), tasa $\delta=2$.
  \item $(0,1)\to(1,1)$: G también falla, tasa $\alpha=1$.
  \item $(1,1)\to(0,1)$: técnico repara G primero, tasa $\beta=4$.
  \item \textbf{$(1,1)\to(1,0)$: NO ocurre} (R no se repara mientras G
        esté fallado).
\end{itemize}

\[
  Q = \begin{pmatrix}
    -1.5 & 1   & 0.5 & 0    \\
    4    & -4.5 & 0   & 0.5  \\
    2    & 0   & -3  & 1    \\
    0    & 0   & 4   & -4
  \end{pmatrix}.
\]

\textbf{Paso 2: resolver $\boldsymbol{\pi}Q = \boldsymbol{0}$.}

Ecuaciones de balance global:

Estado 0: $\;4\pi_1 + 2\pi_2 = 1.5\,\pi_0.$

Estado 1: $\;\pi_0 = 4.5\,\pi_1 - 0.5\,\pi_3.\;$ Reescribimos:
$\pi_0 + 0.5\pi_3 = 4.5\pi_1.$

Estado 2: $\;0.5\pi_0 + 4\pi_3 = 3\pi_2.$

Estado 3: $\;0.5\pi_1 + \pi_2 = 4\pi_3.$

De la ecuación del estado 3: $\pi_3 = (0.5\pi_1+\pi_2)/4$.

Expresemos todo en función de $\pi_0$.  De estado 0 y estado 1,
probemos con balance detallado parcial.  En realidad, al no ser
nacimiento-muerte, resolvamos el sistema directamente.

De (E1): $\pi_0 = 4.5\pi_1 - 0.5\pi_3$.
De (E3): $\pi_3 = 0.125\pi_1 + 0.25\pi_2$.

Sustituimos (E3) en (E1):
$\pi_0 = 4.5\pi_1 - 0.5(0.125\pi_1+0.25\pi_2)
= 4.5\pi_1 - 0.0625\pi_1 - 0.125\pi_2
= 4.4375\pi_1 - 0.125\pi_2.$

De (E2): $0.5\pi_0 + 4(0.125\pi_1+0.25\pi_2) = 3\pi_2$,
es decir $0.5\pi_0 + 0.5\pi_1 + \pi_2 = 3\pi_2$,
así que $0.5\pi_0 + 0.5\pi_1 = 2\pi_2$, lo que da
$\pi_2 = (0.5\pi_0+0.5\pi_1)/2 = 0.25\pi_0+0.25\pi_1$.

Sustituyendo en la expresión de $\pi_0$:
\begin{align*}
  \pi_0 &= 4.4375\pi_1 - 0.125(0.25\pi_0+0.25\pi_1) \\
  &= 4.4375\pi_1 - 0.03125\pi_0 - 0.03125\pi_1 \\
  1.03125\,\pi_0 &= 4.40625\,\pi_1 \\
  \pi_1 &= \frac{1.03125}{4.40625}\,\pi_0 = 0.23404\,\pi_0.
\end{align*}

Entonces:
\begin{align*}
  \pi_2 &= 0.25\pi_0 + 0.25(0.23404\pi_0) = 0.25\pi_0+0.05851\pi_0
  = 0.30851\,\pi_0, \\
  \pi_3 &= 0.125(0.23404\pi_0)+0.25(0.30851\pi_0)
  = 0.02926\pi_0 + 0.07713\pi_0 = 0.10638\,\pi_0.
\end{align*}

Normalización:
\[
  \pi_0(1+0.23404+0.30851+0.10638) = 1
  \quad\Longrightarrow\quad
  1.64894\,\pi_0 = 1
  \quad\Longrightarrow\quad
  \pi_0 \approx 0.6065.
\]
\[
  \boldsymbol{\pi} \approx (0.6065,\;0.1419,\;0.1871,\;0.0645).
\]

\textbf{Verificación:} la suma es $0.6065+0.1419+0.1871+0.0645=1.0000$. \checkmark

\textbf{Paso 3: métricas de estado estable.}

Definimos las variables indicadoras:
$G = \begin{cases}0&\text{si G funciona}\\1&\text{si G fallado}\end{cases}$
y análogamente $R$.  El número de generadores fallados es $N=G+R$.

\emph{Marginales:}
\begin{align*}
  \Prob(G=1) &= \pi_1+\pi_3 = 0.1419+0.0645 = 0.2064, \\
  \Prob(R=1) &= \pi_2+\pi_3 = 0.1871+0.0645 = 0.2516.
\end{align*}

\emph{Valor esperado:}
\[
  \E[N] = \E[G]+\E[R] = 0.2064+0.2516 = 0.4580.
\]

\emph{Varianzas individuales:}
\begin{align*}
  \Var[G] &= 0.2064\times 0.7936 = 0.1638, \\
  \Var[R] &= 0.2516\times 0.7484 = 0.1883.
\end{align*}

\emph{Covarianza:}
\[
  \E[GR] = \Prob(G=1,R=1) = \pi_3 = 0.0645.
\]
\[
  \Cov(G,R) = 0.0645 - 0.2064\times 0.2516 = 0.0645 - 0.0519 = 0.0126.
\]

La covarianza es \emph{positiva}: cuando G falla, el técnico se ocupa de
G, dejando a R sin reparar, lo que aumenta la probabilidad de que R
también esté fallado.  La dependencia estructural (un solo técnico)
genera correlación positiva.

\emph{Varianza de $N$:}
\[
  \Var[N] = \Var[G]+\Var[R]+2\Cov(G,R)
  = 0.1638+0.1883+2(0.0126) = 0.3773.
\]

Si los generadores fueran independientes, la varianza sería
$0.1638+0.1883=0.3521$.  El efecto del técnico compartido agrega
$2\times 0.0126 = 0.0252$ de varianza extra.

\textbf{Paso 4: ¿cuál es la fracción del tiempo sin ningún generador
funcionando?}
\[
  \Prob(\text{apagón total}) = \pi_3 \approx 6.45\%.
\]
\end{ejemplo}

% ======================================================================
%  SECCIÓN 10 — EJEMPLO CON DOS COLAS PEQUEÑAS
% ======================================================================
\section{Ejemplo 5 --- Dos colas acopladas: desbordamiento}

\begin{ejemplo}[Dos cajeros de estacionamiento]
Un estacionamiento tiene dos máquinas expendedoras de tiquetes.  La
máquina~1 (más visible) recibe autos con tasa $\lambda=8$/hora; si está
ocupada, el auto pasa a la máquina~2.  Si ambas están ocupadas, el auto
se va.  La máquina~1 despacha con tasa $\mu_1=10$/hora y la~2 con
$\mu_2=6$/hora.  Cada máquina atiende a un solo auto a la vez (sin cola).

Estados: $(C_1,C_2)\in\{(0,0),(1,0),(0,1),(1,1)\}$, numerados
$0,1,2,3$.

\tcblower

\textbf{Paso 1: generadora.}

\begin{itemize}[nosep]
  \item $(0,0)\to(1,0)$: llega auto, va a máq.~1. Tasa $8$.
  \item $(1,0)\to(0,0)$: máq.~1 termina. Tasa $10$.
  \item $(1,0)\to(1,1)$: llega auto, máq.~1 ocupada, va a máq.~2. Tasa $8$.
  \item $(0,1)\to(1,1)$: llega auto, máq.~1 libre\ldots\ pero espere:
        si máq.~1 está libre el auto va ahí, no a la~2.  Así que:
        $(0,1)\to(1,1)$ con tasa $8$.
  \item $(0,1)\to(0,0)$: máq.~2 termina. Tasa $6$.
  \item $(1,1)\to(0,1)$: máq.~1 termina. Tasa $10$.
  \item $(1,1)\to(1,0)$: máq.~2 termina. Tasa $6$.
\end{itemize}

\[
  Q = \begin{pmatrix}
    -8  & 8   & 0  & 0   \\
    10  & -18 & 0  & 8   \\
    6   & 0   & -14 & 8  \\
    0   & 6   & 10 & -16
  \end{pmatrix}.
\]

\textbf{Paso 2: resolver por balance global.}

(E0): $10\pi_1 + 6\pi_2 = 8\pi_0$.

(E1): $8\pi_0 + 6\pi_3 = 18\pi_1$.

(E2): $10\pi_3 = 14\pi_2 - 6\pi_2 \ldots$ Escribámoslo correctamente.
La ecuación de balance para el estado~2 es: flujo de entrada = flujo de
salida.  Entrada a~2: desde $(0,0)$~no hay transición directa, desde
$(1,1)$: tasa $10\pi_3$.  Salida: $14\pi_2$.  Pero espere,
¿hay transición $(0,0)\to(0,1)$?  No, porque un auto que llega siempre
va primero a máq.~1.  Así que el estado $(0,1)$ solo se alcanza desde
$(1,1)$ cuando máq.~1 termina.

Reescribamos:

(E2): $10\pi_3 = 14\pi_2$.  Entonces $\pi_2 = (10/14)\pi_3 = (5/7)\pi_3$.

(E3): $8\pi_1 + 8\pi_2 = 16\pi_3$.  Espere, revisemos.  Al estado~3
se llega desde~1 (tasa $8$) y desde~2 (tasa $8$).  Se sale con tasa
$16$.  Entonces: $8\pi_1+8\pi_2 = 16\pi_3$, es decir
$\pi_1+\pi_2 = 2\pi_3$.

De (E2): $\pi_2 = 5\pi_3/7$.  Sustituyendo en (E3):
$\pi_1 + 5\pi_3/7 = 2\pi_3$, así $\pi_1 = 2\pi_3-5\pi_3/7 = 9\pi_3/7$.

De (E1): $8\pi_0+6\pi_3 = 18(9\pi_3/7) = 162\pi_3/7$.
Entonces $8\pi_0 = 162\pi_3/7 - 6\pi_3 = (162-42)\pi_3/7 = 120\pi_3/7$.
$\pi_0 = 15\pi_3/7$.

Normalización: $\pi_0+\pi_1+\pi_2+\pi_3 = 1$:
\[
  \frac{15}{7}\pi_3 + \frac{9}{7}\pi_3+\frac{5}{7}\pi_3+\pi_3 = 1
  \quad\Longrightarrow\quad
  \frac{15+9+5+7}{7}\pi_3 = 1
  \quad\Longrightarrow\quad
  \frac{36}{7}\pi_3 = 1.
\]
\[
  \pi_3 = \frac{7}{36}, \quad
  \pi_0 = \frac{15}{36} = \frac{5}{12}, \quad
  \pi_1 = \frac{9}{36} = \frac{1}{4}, \quad
  \pi_2 = \frac{5}{36}.
\]

En decimales:
\[
  \boldsymbol{\pi} \approx (0.4167,\;0.2500,\;0.1389,\;0.1944).
\]
Verificación: $0.4167+0.2500+0.1389+0.1944=1.0000$. \checkmark

\textbf{Paso 3: métricas.}

Número total de máquinas ocupadas $N=C_1+C_2$:
\begin{align*}
  \Prob(N=0) &= \pi_0 = 5/12 \approx 0.4167, \\
  \Prob(N=1) &= \pi_1+\pi_2 = 1/4+5/36 = 14/36 = 7/18 \approx 0.3889, \\
  \Prob(N=2) &= \pi_3 = 7/36 \approx 0.1944.
\end{align*}

\[
  \E[N] = 0\times\frac{5}{12} + 1\times\frac{7}{18} + 2\times\frac{7}{36}
  = \frac{7}{18}+\frac{14}{36} = \frac{14}{36}+\frac{14}{36}
  = \frac{28}{36} = \frac{7}{9} \approx 0.7778.
\]

\[
  \E[N^2] = 0+\frac{7}{18}+4\times\frac{7}{36}
  = \frac{7}{18}+\frac{28}{36} = \frac{14}{36}+\frac{28}{36}
  = \frac{42}{36} = \frac{7}{6} \approx 1.1667.
\]

\[
  \Var[N] = \frac{7}{6}-\Bigl(\frac{7}{9}\Bigr)^2
  = \frac{7}{6}-\frac{49}{81}
  = \frac{7\times 81 - 49\times 6}{6\times 81}
  = \frac{567-294}{486} = \frac{273}{486} = \frac{91}{162}
  \approx 0.5617.
\]

\textbf{Probabilidad de bloqueo y throughput:}
\[
  P_{\text{bloqueo}} = \pi_3 = 7/36 \approx 19.44\%.
\]
\[
  \lambda_{\text{ef}} = 8(1-\pi_3) = 8\times\frac{29}{36}
  = \frac{232}{36} \approx 6.444 \text{ autos/hora}.
\]

\textbf{Análisis por máquina (dos variables):}
\begin{align*}
  \Prob(C_1=1) &= \pi_1+\pi_3 = 1/4+7/36 = 16/36 = 4/9 \approx 0.4444, \\
  \Prob(C_2=1) &= \pi_2+\pi_3 = 5/36+7/36 = 12/36 = 1/3 \approx 0.3333.
\end{align*}

Covarianza:
\[
  \E[C_1 C_2] = \pi_3 = 7/36.
\]
\[
  \Cov(C_1,C_2) = \frac{7}{36}-\frac{4}{9}\times\frac{1}{3}
  = \frac{7}{36}-\frac{4}{27} = \frac{21-16}{108} = \frac{5}{108}
  \approx 0.0463.
\]

La covarianza positiva refleja que cuando la máquina~1 está ocupada, hay
más probabilidad de que los autos desborden a la máquina~2.

Verificación de la varianza:
\begin{align*}
  \Var[C_1] &= \frac{4}{9}\times\frac{5}{9} = \frac{20}{81}, \\
  \Var[C_2] &= \frac{1}{3}\times\frac{2}{3} = \frac{2}{9}, \\
  \Var[N] &= \frac{20}{81}+\frac{2}{9}+2\times\frac{5}{108}
  = \frac{20}{81}+\frac{18}{81}+\frac{10}{108}.
\end{align*}
Llevando a común denominador (324):
$\frac{80}{324}+\frac{72}{324}+\frac{30}{324} = \frac{182}{324}
= \frac{91}{162} \approx 0.5617$. \checkmark
\end{ejemplo}

% ======================================================================
%  SECCIÓN 14 — EJERCICIOS
% ======================================================================
\section{Ejercicios propuestos}

\begin{ejercicio}[Línea de ensamblaje]
Una línea de ensamblaje tiene tres estaciones.  Las piezas llegan a la
estación~1 con tasa $\lambda=12$/hora.  Si la estación~1 está ocupada,
la pieza pasa a la estación~2; si esa también está ocupada, a la~3; si
las tres están ocupadas, la pieza se descarta.  Las tasas de servicio
son $\mu_1=15$, $\mu_2=10$, $\mu_3=8$ piezas/hora.

\begin{enumerate}[label=(\alph*)]
  \item Modele el sistema.
        Simplifique observando que ciertas combinaciones son imposibles
        (¿cuáles y por qué?).
  \item Escriba la matriz $Q$ .
  \item Calcule $\boldsymbol{\pi}$ y la tasa de piezas descartadas.
  \item Calcule $\E[N]$, $\Var[N]$ donde $N$ es el total de estaciones
        ocupadas.
\end{enumerate}
\end{ejercicio}

\begin{ejercicio}[Cola $M/M/1/3$ --- estado estable vs.\ transitorio]
Retome el sistema $M/M/1/3$ del ejercicio de la lectura anterior
($\lambda=10$/hora, $\mu=15$/hora, capacidad 3).

\begin{enumerate}[label=(\alph*)]
  \item Calcule la distribución estacionaria.
  \item Calcule $\E_\pi[N]$ y $\Var_\pi[N]$.
  \item Compare con los valores transitorios que calculó para $t=6$
        minutos.  ¿Se ha acercado el sistema al estado estable?
  \item Calcule la probabilidad de bloqueo y el throughput.
\end{enumerate}
\end{ejercicio}

\begin{ejercicio}[Dos impresoras con desbordamiento]
Una oficina tiene dos impresoras.  Los trabajos llegan con tasa
$\lambda=20$/hora.  Un trabajo va a la impresora~1 si está libre
($\mu_1=25$/hora); si no, va a la impresora~2 ($\mu_2=15$/hora); si
ambas están ocupadas, el trabajo se rechaza.

\begin{enumerate}[label=(\alph*)]
  \item Escriba la generadora $Q$ ($4\times 4$).
  \item Calcule $\boldsymbol{\pi}$.
  \item Determine la fracción de trabajos rechazados.
  \item Calcule $\E[N]$, $\Var[N]$ y $\Cov(C_1,C_2)$.
  \item Un empleado dice: ``como $\lambda<\mu_1+\mu_2$, nunca debería
        haber rechazos''.  Explique por qué está equivocado.
\end{enumerate}
\end{ejercicio}


% ======================================================================
%  BIBLIOGRAFÍA SUGERIDA
% ======================================================================
\section*{Bibliografía sugerida}

\begin{enumerate}[nosep]
  \item W.~J.~Stewart, \textit{Probability, Markov Chains, Queues, and
        Simulation}, Princeton University Press, 2009.
        Capítulos sobre distribuciones estacionarias y balance detallado.
  \item D.~Gross, J.~F.~Shortle, J.~M.~Thompson, C.~M.~Harris,
        \textit{Fundamentals of Queueing Theory}, 4th ed., Wiley, 2008.
        Análisis de estado estable para sistemas de colas.
  \item G.~Bolch, S.~Greiner, H.~de~Meer, K.~S.~Trivedi,
        \textit{Queueing Networks and Markov Chains}, 2nd ed., Wiley, 2006.
        Métodos de solución y métricas de desempeño.
  \item R.~Nelson, \textit{Probability, Stochastic Processes, and Queueing
        Theory}, Springer, 1995.
        Tratamiento accesible del teorema ergódico y aplicaciones.
  \item S.~M.~Ross, \textit{Introduction to Probability Models}, 12th ed.,
        Academic Press, 2019.  Capítulos sobre CTMC y estado estable.
\end{enumerate}

\end{document}